\section{Introduction}
\label{sec:introduction}

Autonomous racing has become a productive research format because it compresses perception, planning, control, and real-time decision-making into a single quantitative task under repeatable constraints~\cite{betz2022autonomous}. In ground racing, the F1TENTH platform standardizes circuits, a vehicle model, metrics, and baselines for continuous control and reinforcement learning~\cite{okelly2020f1tenth}, while full-scale leagues such as the Abu Dhabi Autonomous Racing League (A2RL) push planning and control to the limits of vehicle handling in head-to-head competition~\cite{hoffmann2026a2rl}. In aerial racing, gate-based benchmarks and competitions drove a progression from modular perception--control stacks~\cite{foehn2020alphapilot} to learning-based policies that reach champion-level performance~\cite{kaufmann2023swift}. A common thread is that a well-defined racing benchmark, with fixed geometry, an impartial referee, and machine-readable outcomes, accelerates and disciplines the comparison of autonomy methods.

Underwater robotics has mature simulation tools but lacks a comparable racing benchmark. Existing marine simulators emphasize intervention, inspection, navigation, or sensor realism rather than a compact racing protocol with explicit gates, currents, obstacles, independent scoring, and fleet evaluation. Yet underwater racing is not aerial or ground racing relocated to a different simulator. Vehicle dynamics are slow and strongly damped; global localization is unavailable and must be inferred from drifting inertial and Doppler sensing or from acoustic cues; optical perception is range-limited and degrades with turbidity; acoustic guidance is low-bandwidth and noisy; persistent currents act as a structured disturbance; and physical contact with a gate or another vehicle is ambiguous to detect. These properties change both what a controller can sense and how an evaluator must score behavior.

A benchmark for this domain must therefore separate the autonomy stack from the referee and from the simulator internals. The controller should receive only a documented \emph{official observation}---onboard sensing, an acoustic target cue, and its own race progress---while the referee may use privileged ground-truth state to decide gate crossings, collisions, timeouts, and team scoring. Without this separation, it is impossible to claim that a result was obtained under realistic information, and it is easy to inadvertently leak global pose or exact gate geometry into the control loop.

\mra{} is designed around this separation. It is neither a new hydrodynamic simulator nor a new autonomy algorithm. It is an evaluation layer built on HoloOcean~\cite{potokar2022holoocean} and a BlueROV2-class underwater vehicle~\cite{vonbenzon2022bluerov2}, in which the race---world bounds, ordered gates, start and finish, currents, obstacles, participants, sensors, and referee---is described declaratively and a controller is plugged in by implementing a small interface. The framework supports three usage categories that are kept deliberately distinct so that infrastructure claims are not confused with solved-problem claims:

\begin{itemize}
  \item \textbf{Official benchmark mode}: fixed official tracks, official gate apertures, official observation filtering that excludes ground-truth state, and reproducible referee scoring;
  \item \textbf{Diagnostic mode}: additional logs, participant-state traces, proximity diagnostics, and smoke tests used to validate the benchmark plumbing itself;
  \item \textbf{Experimental mode}: currents, obstacles, current-compensation controllers, and close-proximity fleet interaction, which are available for investigation but are \emph{not} claimed as solved v0.1 results.
\end{itemize}

The contributions of this release-candidate paper are the following.

\begin{enumerate}
  \item A declarative underwater gate-racing benchmark model: world, ordered gates with $1.5\,\text{m}\times1.5\,\text{m}$ apertures, currents, obstacles, participants, sensors, and scoring are specified in a single JSON configuration, with three official clean tracks provided.
  \item An official controller--observation contract and a HoloOcean race runner that keep ground-truth pose, simulator current vectors, and hidden gate geometry out of participant control, together with a controller interface that a custom policy satisfies by implementing three methods.
  \item An independent referee that formalizes gate-crossing geometry, missed gates, collisions, out-of-bounds and stuck conditions, timeouts, ranking, and structured event logs.
  \item A staggered fleet/team mode in which multiple BlueROV2 participants form a single team, maintain independent progress and referee state, and aggregate into one team-level score, with rover--rover proximity diagnostics.
  \item A reproducible v0.1 validation package with a deterministic rule-based baseline, clean single-rover results on three official tracks, a stable staggered two-rover demonstration, and an explicit, honest scoping of current compensation and close-proximity fleet collisions as open work.
\end{enumerate}

The remainder of the paper is organized as follows. Section~\ref{sec:related_work} reviews underwater simulators, autonomous-racing benchmarks, and multi-vehicle underwater cooperation. Section~\ref{sec:system_overview} formalizes the benchmark model and the runtime architecture. Section~\ref{sec:configuration} describes the vehicle, the HoloOcean adapter, the sensor suite, the current model, and the official observation contract. Section~\ref{sec:controller} specifies the controller interface and the baseline controllers. Section~\ref{sec:referee} formalizes gate validation, the referee state machine, and single-rover scoring. Section~\ref{sec:fleet} describes staggered fleet execution and team scoring. Section~\ref{sec:evaluation} reports the experimental evaluation. Section~\ref{sec:discussion} discusses limitations and future work, and Section~\ref{sec:conclusion} concludes.
